\section{Introduction}
% What we do
In this paper, we introduce Grounding DINO 1.5, a series of powerful and practical open-set object detection models developed by IDEA Research. 
% motivation
Object detection is a fundamental task in computer vision, with recent efforts focusing on developing generic detectors capable of performing detection across a wide variety of real-world applications. 
A key strategy for improving model generalization across diverse object categories is the integration of language modality, which has received increasingly more attention and has undergone extensive development in recent research.


% background
Grounding DINO~\cite{GroundingDINO} represents a significant advancement in this area. 
Building on the Transformer-based DINO~\cite{DINO} architecture, it incorporates linguistic information to enable open-set object detection in various scenarios. 
Following GLIP~\cite{GLIP}, Grounding DINO redefines object detection as a phrase grounding task, facilitating large-scale pre-training using both detection and grounding datasets. 
This approach, coupled with self-training on pseudo-labeled grounding data derived from an almost unlimited pool of image-text pairs, enhances the model’s applicability to open-world settings due to its robust architecture and semantic-rich dataset.

% What we do

Building upon the success of Grounding DINO, Grounding DINO 1.5 extends the model's capabilities in two key areas: \emph{stronger} detection performance and \emph{faster} inference speed, designated as the Grounding DINO 1.5 Pro and Grounding DINO 1.5 Edge models, respectively.

The Grounding DINO 1.5 Pro model significantly expands both the model's capacity and the dataset size, with the goal of creating a more potent and versatile open-set object detection model. Specifically, we have upscaled the model by incorporating the pre-trained ViT-L~\cite{EVA-02} architecture and developed a data engine capable of producing over 20 million images with grounding annotations from diverse sources, thereby enriching the model's semantic comprehension. %These enhancements ensure that Grounding DINO 1.5 Pro is more closely aligned with human preferences in visual recognition tasks.

By contrast, the Grounding DINO 1.5 Edge model is tailored for edge devices, focusing on computational efficiency without compromising detection quality. We develop an efficient feature enhancer that leverages only high-level image features, removing the need of multi-scale features. This streamlined approach maintains the model's strong context-aware detection capabilities, after training on the same 20 million images as used for the Pro model.


% results
Extensive results from our experiments validate the superiority of Grounding DINO 1.5.
Specifically, Grounding DINO 1.5 Pro achieves a 54.3 AP on the COCO detection zero-shot transfer benchmark and simultaneously achieves a 55.7 AP and a 47.6 AP on the LVIS-minival and LVIS-val zero-shot transfer benchmarks, respectively, setting new records on these benchmarks. Moreover, under TensorRT optimization, the Grounding DINO 1.5 Edge model reaches a speed of 75.2 FPS and achieves a zero-shot performance of 36.2 AP on the LVIS-minival benchmark, making it more suitable for edge computing scenarios.

